\chapter{Unbiased Estimation}

This chapter is devoted to estimators that are \emph{unbiased}, that is, estimators
whose expectation equals the target parameter.We will develop a
classical optimality theory, including criteria for comparing and improving
unbiased estimators.

The material follows the classical theory of point estimation, as presented for
example in Lehmann and Casella (1998, Chapter~2). 
%Related topics, such as
%equiv arrant estimation and large-sample (asymptotic) theory, are treated in later
%chapters of these notes.

\section{Unbiased Estimators}%
\label{sec:ube}

\paragraph{Setting}
We consider the following framework:
\begin{itemize}
    \item Observation: \(X\), taking values in a sample space \(\cc{X}\).
    \item Statistical model: \(\cc{P} = \{\mathsf{P}_\theta : \theta \in \Theta\}\).
    \item Parameter of interest: \(\gamma(\theta) \in \R\), a real-valued function of \(\theta\).
\end{itemize}

\begin{definition}[]
    An estimator \(T = T(X)\) is called \textbf{unbiased} for \(\gamma(\theta)\) if
    \[
        \mathsf{E}_\theta[T] = \gamma(\theta) \quad \text{for all } \theta \in \Theta,
    \]
    where the expectation is taken with respect to \(X \sim \mathsf{P}_\theta\).

    If an unbiased estimator exists for \(\gamma(\theta)\), the parameter \(\gamma(\theta)\) is called \textbf{U--estimable}.
\end{definition}

\myparagraph{Motivation and Intuition}
Unbiased estimators are attractive because, on average, they return the correct value of the target parameter. That is, repeated applications of the estimator will not systematically overestimate or underestimate the quantity of interest.  

In terms of mean squared error (MSE),
\[
\text{MSE}(T) = \bigl(\text{Bias}(T)\bigr)^2 + \mathsf{Var}_\theta(T),
\]
unbiased estimators have zero bias, so minimizing MSE reduces to minimizing variance. Therefore, when restricting attention to unbiased estimators, the natural question is:

\emph{Does there exist an unbiased estimator with minimal variance?}

\subsection*{Discussion}
The notion of unbiasedness depends on the underlying model \(\cc{P}\) and the parameter of interest \(\gamma(\theta)\). For example:
\begin{itemize}
    \item \(X\) may be a single observation, a vector of independent draws, or even a matrix of observations.
    \item \(\gamma(\theta)\) could be a component of a multivariate parameter, such as the mean of a Gaussian distribution.
\end{itemize}

Unbiased estimators are often preferred when constructing confidence intervals, because the usual $\pm$ margin-of-error construction is centered around the estimator. A biased estimator would systematically shift this interval, potentially distorting coverage probabilities.

Finally, note that not every parameter admits an unbiased estimator. 
%When an unbiased estimator exists, we say the parameter is U--estimable. 
Restricting attention to unbiased estimators removes pathological cases (such as constant estimators that trivially minimize MSE but are useless for estimating a varying parameter), allowing the meaningful study of \emph{optimal unbiased estimators}, i.e., those with minimal variance among all unbiased estimators.


\section{UMVUE}%
\label{sec:UMVUE}
\begin{definition}[]
    Let \(T\) be an unbiased estimator of \(\gamma(\theta)\). If for any other unbiased estimator \(T'\) of \(\gamma(\theta)\), it holds that 
    \[
        \mathsf{Var}_{\theta}[T] \le \mathsf{Var}_{\theta}[T']\quad \forall \theta \in \Theta,
    \]
    then \(T\) is called \textbf{UMVUE} (\emph{Uniform Minimum Variance Unbiased Estimator}) of \(\gamma(\theta)\).

That is, among all unbiased estimators of \( \gamma(\theta) \), the estimator \( T \) achieves the
smallest variance uniformly over the entire parameter space.
\end{definition}

\begin{pro}[Uniqueness of the UMVUE]
\label{prop:UMVUE-unique}
Let \( T \) be a UMVUE of \( \gamma(\theta) \) such that
\[
\mathsf{Var}_{\theta}(T) < \infty, \quad \forall \theta \in \Theta.
\]
Then \( T \) is unique in the following sense: if \( T' \) is another UMVUE of
\( \gamma(\theta) \), then
\[
\mathsf{P}_{\theta}(T = T') = 1, \quad \forall \theta \in \Theta.
\]
\end{pro}

\begin{proof}
Assume that \( T' \) is another UMVUE of \( \gamma(\theta) \).
Since both \( T \) and \( T' \) are unbiased, their average
\[
T^{\ast} := \frac{1}{2}(T + T')
\]
is also an unbiased estimator of \( \gamma(\theta) \).

Because \( T \) is a UMVUE, its variance cannot exceed that of any other unbiased
estimator. In particular,
\[
\mathsf{Var}_{\theta}(T) \le \mathsf{Var}_{\theta}(T^{\ast}),
\quad \forall \theta \in \Theta.
\]
A direct computation yields
\[
\mathsf{Var}_{\theta}(T^{\ast})
= \frac{1}{4}\mathsf{Var}_{\theta}(T)
+ \frac{1}{4}\mathsf{Var}_{\theta}(T')
+ \frac{1}{2}\mathsf{Cov}_{\theta}(T, T').
\]
Since both \( T \) and \( T' \) are UMVUEs, they must have the same variance for
every \( \theta \). Applying the Cauchy--Schwarz inequality,
\[
\mathsf{Cov}_{\theta}(T, T')
\le \sqrt{\mathsf{Var}_{\theta}(T)\mathsf{Var}_{\theta}(T')},
\]
we obtain
\[
\mathsf{Var}_{\theta}(T^{\ast})
\le \mathsf{Var}_{\theta}(T).
\]
Combining this with the earlier inequality shows that equality must hold
throughout. In particular, equality in the Cauchy--Schwarz inequality implies
\[
\operatorname{Corr}_{\theta}(T, T') = 1,
\]
and therefore there exist constants \( a, b \in \mathbb{R} \) such that
\[
\mathsf{P}_{\theta}(T' = aT + b) = 1, \quad \forall \theta \in \Theta.
\]

Because \( T \) and \( T' \) are unbiased for the same target,
\( \mathsf{E}_{\theta}[T] = \mathsf{E}_{\theta}[T'] \) implies \( b = 0 \).
Moreover, equality of variances yields
\[
\mathsf{Var}_{\theta}(T')
= a^{2}\mathsf{Var}_{\theta}(T)
= \mathsf{Var}_{\theta}(T),
\]
so \( a = 1 \). Hence,
\[
\mathsf{P}_{\theta}(T = T') = 1, \quad \forall \theta \in \Theta,
\]
which proves uniqueness.
\end{proof}

\begin{ex}[Binomial model]
Let \( X \sim \operatorname{Binomial}(n,\theta) \), where \( n \) is known and
\( \theta \in (0,1) \).

For any estimator \( T = T(X) \), its expectation is given by
\[
\mathsf{E}_{\theta}[T]
= \sum_{x=0}^{n} T(x)\binom{n}{x}\theta^{x}(1-\theta)^{n-x}.
\]
Since the binomial probability mass function is a polynomial in \( \theta \) of
degree at most \( n \), the expectation \( \mathsf{E}_{\theta}[T] \) is itself a
polynomial in \( \theta \) of degree at most \( n \).
Moreover, by choosing the values \( T(0),\dots,T(n) \) appropriately, any
polynomial of degree at most \( n \) can be represented in this way.

Consequently, the parameters \( \gamma(\theta) \) that are unbiasedly estimable
in the binomial model are precisely the polynomials in \( \theta \) of degree
at most \( n \).
Thus, non-polynomial functions such as
\( \gamma(\theta) = \sqrt{\theta} \) are not unbiasedly estimable.

As a simple example, the parameter \( \gamma(\theta) = \theta \) is unbiasedly
estimable. The estimator
\[
\widehat{\theta} = \frac{X}{n}
\]
is unbiased for \( \theta \) and, in fact, is the UMVUE of \( \theta \) in the
binomial model; see Lehmann and Casella (1998, Example~2.3.1).
\end{ex}

\begin{ex}[Gaussian model]
Let \(X_1,\dots,X_n\) be i.i.d.\ \(\cc{N}(\mu,\sigma^2)\), with
\(\theta=(\mu,\sigma^2)\in\Theta=\R\times(0,\infty)\).

In this model, several UMVUEs can be identified explicitly.
A later theorem shows that the
\begin{itemize}
    \item Sample mean
    \[
        \ol X_n = \frac{1}{n}\sum\limits_{i=1}^{n} X_i
    \]
    is the UMVUE of \(\mu\);
    \item Sample variance
    \[
        s^2 = \frac{1}{n-1}\sum\limits_{i=1}^{n}(X_i-\ol X_n)^2
    \]
    is the UMVUE of \(\sigma^2\).
\end{itemize}

If the target parameter is changed from \(\sigma^2\) to \(\sigma\), unbiasedness is no longer preserved
by taking square roots: \(s=\sqrt{s^2}\) is biased for \(\sigma\).
Nevertheless, there exists an unbiased estimator of \(\sigma\), given by
\[
    2^{1/2}\frac{\Gamma\!\left(\frac{n-1}{2}\right)}{\Gamma\!\left(\frac{n}{2}\right)}
    \cdot
    \sqrt{\sum\limits_{i=1}^{n}(X_i-\ol X_n)^2},
\]
which is the UMVUE of \(\sigma\).
\end{ex}

\begin{note}
In this Gaussian example, the maximum likelihood estimator
\[
    \hat\sigma^2 = \frac{1}{n}\sum\limits_{i=1}^{n}(X_i-\ol X_n)^2
\]
has smaller mean square error than the UMVUE \(s^2\), despite being biased.
More generally, among estimators of the form
\[
    \frac{1}{a}\sum\limits_{i=1}^{n}(X_i-\ol X_n)^2,
\]
the choice \(a=n+1\) minimizes the mean square error.
This highlights that UMVUE optimality depends on the unbiasedness criterion and
does not necessarily align with MSE optimality.
\end{note}

\subsection{Behrens--Fisher Problem}%
\label{sub:bfp}

\begin{pro}[Behrens--Fisher Problem]
Consider two independent samples \(X_1,\dots ,X_{n}\) and \(Y_1,\dots ,Y_{m}\), where
\[
X_i \stackrel{\text{i.i.d.}}{\sim} \cc{N}(\mu,\sigma^2),
\qquad
Y_j \stackrel{\text{i.i.d.}}{\sim} \cc{N}(\mu,\tau^2),
\]
with unknown parameter \(\theta=(\mu,\sigma^2,\tau^2)\).
Although unbiased estimators of the common mean \(\mu\) exist, there is no UMVUE of \(\mu\).
\end{pro}

\begin{proof}

Suppose, for contradiction, that there exists a UMVUE \(\hat{\mu}^{\ast}\) for \(\mu\) in the full model.

To analyze optimality, consider a submodel in which the variance ratio
\[
\frac{\tau^2}{\sigma^2} = a
\]
is fixed and known, with \(a>0\). In this submodel, the optimal way to combine information from the two samples is determined by the relative variances. It can be shown that the UMVUE of \(\mu\) in this submodel is
\[
\hat{\mu}_{a}
=
\frac{\sum\limits_{i=1}^{n} X_{i}
+
\frac{1}{a}\sum\limits_{j=1}^{m} Y_{j}}
{n+\frac{1}{a}m}.
\]
This estimator assigns more weight to the sample with smaller variance and is unbiased for \(\mu\).

Importantly, \(\hat{\mu}_{a}\) remains an unbiased estimator of \(\mu\) even in the full model with arbitrary \((\mu,\sigma^2,\tau^2)\). Since \(\hat{\mu}^{\ast}\) is assumed to be a UMVUE in the full model, it must also be a UMVUE within every submodel. By the uniqueness of UMVUEs, this implies
\[
\hat{\mu}^{\ast} = \hat{\mu}_{a}
\quad \text{Lebesgue--almost everywhere}.
\footnote{Null sets of any multivariate normal distribution coincide with sets of Lebesgue measure zero.}
\]

However, the estimator \(\hat{\mu}_{a}\) depends explicitly on the value of \(a\). Since \(a>0\) was arbitrary, the above almost sure equality would have to hold simultaneously for all \(a>0\), which is impossible. This contradiction shows that no UMVUE for \(\mu\) exists in the Behrens--Fisher problem.

The underlying intuition is that an optimal combination of the sample means must weight them according to their relative variances. If the variance ratio were known, the optimal weighting would be clear. When the ratio is unknown, no single unbiased estimator can uniformly dominate all variance--specific combinations, preventing the existence of a UMVUE.
\end{proof}

\subsection{\(L_{2}\)--Orthogonality}%
\label{sub:l2}

A useful way to understand optimal unbiased estimation is through the geometry of
\(L_{2}\) spaces.  Unbiased estimators of a fixed target \(\gamma(\theta)\) form an
affine space, and differences of such estimators lie in a linear space of estimators
with mean zero.  Since variance is simply the squared \(L_{2}\)-norm, the problem of
finding a UMVUE becomes a problem of
orthogonal projection.  The following theorem formalizes this idea.

\begin{theo}[Characterization via \(L_{2}\)-orthogonality]
\label{theo:l2}
Let \(T\) be an unbiased estimator of \(\gamma(\theta)\) with
\(\mathsf{Var}_{\theta}[T] < \infty\) for all \(\theta \in \Theta\).
For each \(\theta \in \Theta\), define
\[
    \cc{U}(\theta)
    =
    \left\{
        U \::\:
        \mathsf{E}_{\theta'}[U] = 0 \ \forall \theta' \in \Theta,
        \ \mathsf{Var}_{\theta}[U] < \infty
    \right\}.
\]
Then \(T\) is a UMVUE of \(\gamma(\theta)\) if and only if
\[
    \mathsf{Cov}_{\theta}[T,U]
    =
    \mathsf{E}_{\theta}[T \cdot U]
    =
    0
    \quad
    \forall \theta \in \Theta,
    \ \forall U \in \cc{U}(\theta).
\]
\end{theo}

\medskip

The condition above has a natural geometric interpretation.
For \(U \in \cc{U}(\theta)\), we have \(\mathsf{E}_{\theta}[U]=0\), and therefore
\[
    \mathsf{Cov}_{\theta}[T,U]
    =
    \mathsf{E}_{\theta}\!\left[(T-\mathsf{E}_{\theta}[T])U\right]
    =
    \mathsf{E}_{\theta}[T \cdot U].
\]
Thus the covariance plays the role of an inner product in \(L_{2}(\mathsf{P}_{\theta})\).
The theorem states that a UMVUE is precisely an unbiased estimator that is orthogonal,
at every \(\theta\), to the subspace of unbiased estimators of zero.
As in Euclidean geometry, orthogonality characterizes the point of minimal distance,
here corresponding to minimal variance.

\begin{remark}[Uniqueness of the UMVUE]
The orthogonality characterization immediately implies uniqueness.
If both \(T\) and \(T'\) are UMVUEs of \(\gamma(\theta)\), then
\(T-T' \in \cc{U}(\theta)\), and hence
\begin{align*}
    \mathsf{Var}_{\theta}[T-T']
    &=
    \mathsf{Cov}_{\theta}[T-T',T-T']  \\
    &=
    \mathsf{Cov}_{\theta}[T,T-T']
    -
    \mathsf{Cov}_{\theta}[T',T-T']
    =
    0.
\end{align*}
Consequently, \(\mathsf{P}_{\theta}(T=T')=1\) for all \(\theta\in\Theta\).
\end{remark}

\begin{proof}
\leavevmode
\begin{description}
\item[\(\Longrightarrow)\)]
Assume that \(T\) satisfies the orthogonality condition.
Let \(T'\) be any other unbiased estimator of \(\gamma(\theta)\), so that
\(\mathsf{E}_{\theta}[T-T']=0\) for all \(\theta\in\Theta\).

Fix \(\theta\in\Theta\).
If \(\mathsf{Var}_{\theta}[T']=\infty\), then trivially
\(\mathsf{Var}_{\theta}[T]\le \mathsf{Var}_{\theta}[T']\).
Otherwise, suppose \(\mathsf{Var}_{\theta}[T']<\infty\).
Since both \(T\) and \(T'\) have finite variance under \(\mathsf{P}_{\theta}\),
their difference satisfies
\[
    \mathsf{Var}_{\theta}[T-T']
    =
    \mathsf{Var}_{\theta}[T]
    +
    \mathsf{Var}_{\theta}[T']
    -
    2\mathsf{Cov}_{\theta}[T,T']
    < \infty,
\]
and hence \(T-T' \in \cc{U}(\theta)\).

By orthogonality,
\[
    0
    =
    \mathsf{Cov}_{\theta}[T,T-T']
    =
    \mathsf{Var}_{\theta}[T]
    -
    \mathsf{Cov}_{\theta}[T,T'].
\]
Therefore,
\[
    \mathsf{Var}_{\theta}[T]
    =
    \mathsf{Cov}_{\theta}[T,T']
    \le
    \sqrt{\mathsf{Var}_{\theta}[T]\,
          \mathsf{Var}_{\theta}[T']}
\]
by the Cauchy--Schwarz inequality.
Dividing by \(\sqrt{\mathsf{Var}_{\theta}[T]}\) yields
\(
\mathsf{Var}_{\theta}[T]
\le
\mathsf{Var}_{\theta}[T'].
\)
Since \(T'\) was arbitrary, \(T\) is a UMVUE.

\item[\(\Longleftarrow)\)]
Suppose \(T\) is a UMVUE and let \(U \in \cc{U}(\theta)\).
For any \(a \in \mathbb{R}\), the estimator \(T+aU\) is unbiased for
\(\gamma(\theta)\).
Optimality of \(T\) implies that the function
\(a \mapsto \mathsf{Var}_{\theta}[T+aU]\) is minimized at \(a=0\),
which forces
\(
\mathsf{Cov}_{\theta}[T,U]=0.
\)
\end{description}
\end{proof}

This characterization is often useful in practice.
In some models, such as the binomial family, one can directly verify the orthogonality
condition to identify a UMVUE.
Conversely, the failure of such an orthogonality relation may be used to show that no
UMVUE exists.
In this sense, the theorem provides both a constructive and a diagnostic tool in the
theory of unbiased estimation.

\subsubsection*{Some Preparation on Conditioning}%
\label{ssub:spc}

Conditioning provides a powerful and often simpler alternative to direct
\(L_{2}\)-orthogonality arguments when constructing good estimators.
Beyond its interpretation as an averaging or integration operation,
conditional expectation can be understood as the optimal prediction of one
random variable based on the information contained in another.
This viewpoint naturally leads to variance decompositions and explains why
conditioning typically reduces variability.

Throughout this section, let \(Y\) and \(Z\) be random variables such that
\(\mathsf{Var}[Y] < \infty\).

\begin{lemma}[Basic identities for conditioning]
\label{lem:conditioning}
The following properties hold:
\begin{enumerate}[label=\roman*)]
    \item \emph{Law of total expectation (tower rule):}
    \[
        \mathsf{E}[Y] = \mathsf{E}[\mathsf{E}[Y \mid Z]].
    \]

    \item \emph{Law of total variance:}
    \[
        \mathsf{Var}[Y]
        =
        \mathsf{Var}[\mathsf{E}[Y \mid Z]]
        +
        \mathsf{E}[\mathsf{Var}[Y \mid Z]].
    \]
    \label{ite:conditioning2}

    \item \emph{Degeneracy conditions:}
    \[
        \mathsf{Var}[\mathsf{E}[Y \mid Z]] = 0
        \iff
        \mathsf{E}[Y \mid Z] \ \text{is a.s.\ constant},
    \]
    and
    \[
        \mathsf{E}[\mathsf{Var}[Y \mid Z]] = 0
        \iff
        Y = \mathsf{E}[Y \mid Z] \quad \text{a.s.}
    \]
    \label{ite:conditioning3}
\end{enumerate}
\end{lemma}

\medskip

These identities reveal how the variability of \(Y\) decomposes into two
distinct components.
The term \(\mathsf{Var}[\mathsf{E}[Y \mid Z]]\) measures the variability explained
by the information in \(Z\), while \(\mathsf{E}[\mathsf{Var}[Y \mid Z]]\) captures
the remaining randomness after conditioning.
In particular, conditioning can never increase variance.

\begin{proof}
\leavevmode
\begin{enumerate}[label=\roman*)]
\item
The first identity is the law of total expectation.

\item
Observe that
\begin{align*}
    \mathsf{Var}\!\left[\mathsf{E}[Y \mid Z]\right]
    &=
    \mathsf{E}\!\left[\left(\mathsf{E}[Y \mid Z]\right)^2\right]
    -
    \left(\mathsf{E}[\mathsf{E}[Y \mid Z]]\right)^2 \\
    &=
    \mathsf{E}\!\left[\left(\mathsf{E}[Y \mid Z]\right)^2\right]
    -
    \left(\mathsf{E}[Y]\right)^2,
\end{align*}
while
\begin{align*}
    \mathsf{E}\!\left[\mathsf{Var}[Y \mid Z]\right]
    &=
    \mathsf{E}\!\left[
        \mathsf{E}[Y^2 \mid Z]
        -
        \left(\mathsf{E}[Y \mid Z]\right)^2
    \right] \\
    &=
    \mathsf{E}[Y^2]
    -
    \mathsf{E}\!\left[\left(\mathsf{E}[Y \mid Z]\right)^2\right].
\end{align*}
Adding these two expressions yields
\[
    \mathsf{Var}[\mathsf{E}[Y \mid Z]]
    +
    \mathsf{E}[\mathsf{Var}[Y \mid Z]]
    =
    \mathsf{E}[Y^2]
    -
    \left(\mathsf{E}[Y]\right)^2
    =
    \mathsf{Var}[Y].
\]

\item
Note that
\begin{align*}
    \mathsf{E}\!\left[\left(Y - \mathsf{E}[Y \mid Z]\right)^2\right]
    &=
    \mathsf{E}[Y^2]
    -
    2\mathsf{E}\!\left[Y\,\mathsf{E}[Y \mid Z]\right]
    +
    \mathsf{E}\!\left[\left(\mathsf{E}[Y \mid Z]\right)^2\right] \\
    &=
    \mathsf{E}[Y^2]
    -
    2\mathsf{E}\!\left[\mathsf{E}\!\left[Y\,\mathsf{E}[Y \mid Z] \mid Z\right]\right]
    +
    \mathsf{E}\!\left[\left(\mathsf{E}[Y \mid Z]\right)^2\right] \\
    &=
    \mathsf{E}[Y^2]
    -
    \mathsf{E}\!\left[\left(\mathsf{E}[Y \mid Z]\right)^2\right] \\
    &=
    \mathsf{E}[\mathsf{Var}[Y \mid Z]].
\end{align*}
Thus \(\mathsf{E}[\mathsf{Var}[Y \mid Z]] = 0\) if and only if
\(Y = \mathsf{E}[Y \mid Z]\) almost surely.
Similarly, \(\mathsf{Var}[\mathsf{E}[Y \mid Z]] = 0\) holds if and only if
\(\mathsf{E}[Y \mid Z]\) is almost surely constant.
\end{enumerate}
\end{proof}

From a geometric perspective, \(\mathsf{E}[Y \mid Z]\) is the best approximation to
\(Y\) among all functions of \(Z\) in the \(L_{2}\)-sense.
That is, it minimizes
\[
    \mathsf{E}\!\left[(Y - g(Z))^2\right]
\]
over all measurable functions \(g\).
This variance-minimization property explains why conditioning is such an effective
tool in estimation theory: replacing a statistic by its conditional expectation
given additional information can only reduce variance.

\subsection{Rao--Blackwell Theorem}%
\label{sub:rb}

A central principle in estimation theory is that \emph{conditioning reduces variance}.
The Rao--Blackwell theorem formalizes this idea and shows how an arbitrary unbiased
estimator can be systematically improved by conditioning on a sufficient statistic.
This result provides one of the main practical tools for constructing optimal
unbiased estimators.

\begin{theo}[Rao--Blackwell]
\label{theo:rb}
Let \(T\) be an unbiased estimator of \(\gamma(\theta)\) with
\(\mathsf{Var}_{\theta}[T] < \infty\) for all \(\theta \in \Theta\).
Suppose that \(S\) is a sufficient statistic for \(\theta\).
Define
\[
    T^{*}
    :=
    \mathsf{E}_{*}[T \mid S]
    \equiv
    \mathsf{E}_{\theta}[T \mid S].
\]
Then \(T^{*}\) is an unbiased estimator of \(\gamma(\theta)\) and satisfies
\[
    \mathsf{Var}_{\theta}[T^{*}]
    \le
    \mathsf{Var}_{\theta}[T]
    \quad
    \forall \theta \in \Theta.
\]
Moreover,
\[
    \mathsf{Var}_{\theta}[T^{*}]
    <
    \mathsf{Var}_{\theta}[T]
    \quad
    \text{unless }
    T = T^{*} \ \text{a.s. under } \mathsf{P}_{\theta}.
\]
\end{theo}
\medskip
\begin{remark}[]
    
The logic of the theorem has two complementary components.
First, conditioning on any statistic can only decrease variance, since
\[
    \mathsf{Var}_{\theta}[T]
    =
    \mathsf{Var}_{\theta}[\mathsf{E}_{\theta}[T \mid S]]
    +
    \mathsf{E}_{\theta}[\mathsf{Var}_{\theta}[T \mid S]],
\]
and the second term is nonnegative.
Second, sufficiency of \(S\) ensures that the conditional expectation
\(\mathsf{E}_{*}\left[T|S\right] \equiv \mathsf{E}_{\theta}[T \mid S]\) does not depend on the unknown parameter
\(\theta\), so that \(T^{*}\) is a \textbf{ well-defined statistic}, i.e.\ a function of the data alone.
\end{remark}

\begin{proof}
Unbiasedness follows from the law of total expectation:
\[
    \mathsf{E}_{\theta}[T^{*}]
    =
    \mathsf{E}_{\theta}[\mathsf{E}_{*}[T \mid S]]
    =
    \mathsf{E}_{\theta}[\mathsf{E}_{\theta}[T \mid S]]
    =
    \mathsf{E}_{\theta}[T]
    =
    \gamma(\theta),
    \quad
    \forall \theta \in \Theta.
\]
To compare variances, apply the law of total variance by part (\ref{ite:conditioning2} of
\autoref{lem:conditioning}:
\[
    \mathsf{Var}_{\theta}[T]
    =
    \mathsf{Var}_{\theta}[T^{*}]
    +
    \mathsf{E}_{\theta}\![
        \underbrace{\mathsf{Var}_{*}[T \mid S]}_{\ge 0}
    ]
    \ge
    \mathsf{Var}_{\theta}[T^{*}].
\]
Equality holds if and only if
\(\mathsf{E}_{\theta}[\mathsf{Var}_{*}[T \mid S]] = 0\),
which by (\ref{ite:conditioning3} of \autoref{lem:conditioning} occurs precisely when
\(T = \mathsf{E}_{\theta}[T \mid S]\) almost surely.
\end{proof}

\medskip

The Rao--Blackwell theorem thus provides a universal variance-improvement
procedure: starting from any unbiased estimator, conditioning on a sufficient
statistic yields a new estimator that is at least as good, and typically strictly
better.
This result will serve as the foundation for constructing uniformly minimum
variance unbiased estimators in subsequent sections.

\subsection{Complete Statistics}%
\label{sub:cs}

Completeness is an additional structural property of statistics that plays a
central role in identifying estimators that cannot be further improved by
conditioning.
As with sufficiency, completeness is always defined relative to a fixed
statistical model.

\medskip

\begin{notation}[]
For a model \(\cc{P}=\{\mathsf{P}_{\theta} : \theta \in \Theta\}\) and random
variables \(X\) and \(Y\), we write
\[
    X = Y \;[\cc{P}\text{--a.e.}]
\]
if
\(
\mathsf{P}_{\theta}(X=Y)=1
\)
for all \(\theta \in \Theta\).
That is, the equality holds almost surely under every distribution in the model.
\end{notation}

\medskip

\begin{definition}[Complete statistic]
A statistic \(S\) is said to be \textbf{complete} for the model
\(\cc{P}=\{\mathsf{P}_{\theta} : \theta \in \Theta\}\) if for every measurable
function \(h\),
\[
    \mathsf{E}_{\theta}[h(S)] = 0
    \quad
    \forall \theta \in \Theta
    \;\Longrightarrow\;
    h(S) = 0 \;[\cc{P}\text{--a.e.}].
\]
\end{definition}

\medskip

Thus, completeness requires that the only function of \(S\) whose expectation
vanishes under all distributions in the model is the trivial one.
In particular, while a random variable may have expectation zero under many
different distributions, completeness rules out such nontrivial cancellations
when the random variable is a function of a complete statistic.

The importance of completeness becomes apparent in conjunction with the
Rao--Blackwell theorem.
Conditioning an unbiased estimator on a sufficient statistic can only reduce
variance; however, unless the sufficient statistic is complete, further
variance reduction may still be possible.
Complete sufficient statistics are precisely those for which this process
cannot be continued, leading to estimators that are optimal among all unbiased
estimators.

\medskip

As a structural result, known as \textbf{Bahadur's theorem},
\[
    \text{complete and sufficient}
    \;\Longrightarrow\;
    \text{minimal sufficient}.
\]
The converse does not hold in general: a minimal sufficient statistic need not
be complete.

\medskip

Several examples of complete statistics will be discussed in the following
sections.

\subsection{Lehmann--Scheff\'e Theorem}%
\label{sub:ls}

The Lehmann--Scheff\'e theorem formalizes the idea that conditioning on a statistic that is both sufficient and complete does not merely improve an unbiased estimator, but actually produces the \emph{unique optimal} one.

\begin{theo}[Lehmann--Scheff\'e]
\label{theo:ls}
In the model \(\cc{P}=\{\mathsf{P}_{\theta} : \theta\in \Theta\}\), let \(T\) be an unbiased estimator of \(\gamma(\theta)\) with \(\mathsf{Var}_{\theta}(T)<\infty\) for all \(\theta\in\Theta\).  
If \(S\) is a sufficient and complete statistic, then
\[
T^{*} \coloneqq \mathsf{E}_{*}[T \mid S]
\]
is the UMVUE of \(\gamma(\theta)\).
\end{theo}

\myparagraph{Motivation and intuition.}
Conditioning on a sufficient statistic \(S\) (via Rao--Blackwell) always improves variance among unbiased estimators.  
However, conditioning alone does not guarantee optimality: if we condition on a merely sufficient (or even minimal sufficient) statistic, different unbiased estimators may still lead to different conditional expectations.

Completeness is the additional assumption that rules this out.  
It says that any function of \(S\) whose expectation is zero for all \(\theta\) must be the zero function.  
This forces all unbiased estimators that are functions of \(S\) to coincide almost surely, leaving only a single candidate.  
That candidate must therefore be the uniformly minimum variance unbiased estimator.

\begin{proof}[]
As in the proof of the Rao--Blackwell \autoref{theo:rb}, the estimator \(T^{*}\) is unbiased because
\[
\mathsf{E}_{\theta}[T^{*}]
= \mathsf{E}_{\theta}\!\left[\mathsf{E}_{*}[T\mid S]\right]
= \mathsf{E}_{\theta}\!\left[\mathsf{E}_{\theta}[T\mid S]\right]
= \mathsf{E}_{\theta}[T]
= \gamma(\theta).
\]

Now let \(T'\) be any other unbiased estimator of \(\gamma(\theta)\).  
By sufficiency of \(S\), using Rao--Blackwell yields an improved unbiased estimator
\[
T'' \coloneqq \mathsf{E}_{*}[T' \mid S],
\]
which is a function of \(S\).

Both \(T^{*}\) and \(T''\) are functions of \(S\) and satisfy
\[
\mathsf{E}_{\theta}[T^{*}-T''] = \gamma(\theta)-\gamma(\theta)=0
\quad \text{for all } \theta\in\Theta.
\]
Hence, \(T^{*}-T''\) is a function of \(S\) with expectation zero for all \(\theta\).
By completeness of \(S\), this implies
\[
T^{*} = T'' \qquad [\cc{P}\text{-a.e.}].
\]

Therefore, every unbiased estimator, after conditioning on \(S\), collapses to the same estimator \(T^{*}\).  
This estimator is thus the \(\cc{P}\)-almost surely unique UMVUE of \(\gamma(\theta)\).
\end{proof}

\begin{note}
If an unbiased estimator \(T\) is already a function of the complete and sufficient statistic \(S\), then it is automatically the UMVUE, since
\[
T = \mathsf{E}_{*}[T\mid S].
\]
\end{note}

\begin{ex}[Binomial model]
    \label{ex:bin}
Let \(X \sim \mathrm{Binomial}(n,\theta)\) with \(\theta\in(0,1)\).
Since the entire sample reduces to the total number of successes, \(X\) is (trivially) a sufficient statistic.

To verify completeness, let \(h\) be any function such that
\[
\mathsf{E}_{\theta}[h(X)] = 0
\quad \text{for all } \theta\in(0,1).
\]
Then
\[
\mathsf{E}_{\theta}[h(X)]
= \sum_{x=0}^{n} h(x)\binom{n}{x}\theta^{x}(1-\theta)^{n-x},
\]
which is a polynomial in \(\theta\) of degree at most \(n\).
If this polynomial vanishes for all \(\theta\in(0,1)\), then it must be the zero polynomial, and hence all of its coefficients are zero.

In particular, the constant term satisfies
\[
h(0)\binom{n}{0} = h(0) = 0.
\]
Proceeding inductively, one concludes that \(h(x)=0\) for all \(x=0,\dots,n\).
Therefore, \(X\) is a complete statistic.

We conclude that \(X\) is both sufficient and complete.

Now consider estimation. Since \(\mathsf{E}_{\theta}[X]=n\theta\), the estimator
\[
\hat{\theta} = \frac{1}{n}X
\]
is unbiased for \(\theta\).
As a function of the complete and sufficient statistic \(X\), it follows from the Lehmann--Scheff\'e theorem that \(\hat{\theta}\) is the UMVUE of \(\theta\).

Similarly, to estimate \(\theta^{2}\), note that
\[
\mathsf{E}_{\theta}\!\left[\frac{X(X-1)}{n(n-1)}\right] = \theta^{2}.
\]
Thus,
\[
\widehat{\theta^{2}} = \frac{X(X-1)}{n(n-1)}
\]
is an unbiased estimator and, being a function of the complete and sufficient statistic \(X\), is the UMVUE of \(\theta^{2}\).
\end{ex}

\begin{ex}[Uniform distributions]
    \label{ex:unif}
Let \(X_{1},\dots,X_{n}\) be i.i.d.\ Uniform\((0,\theta)\) with \(\theta>0\).
The parameter \(\theta\) is the unknown upper endpoint of the support.

As seen earlier, the statistic
\[
T(X) = \max\{X_{1},\dots,X_{n}\}
\]
is sufficient. Intuitively, all information about the endpoint \(\theta\) is contained in the largest observation.

To determine the distribution of \(T\), we compute its distribution function. For \(0\le t\le\theta\),
\[
\mathsf{P}_{\theta}(T\le t)
= \mathsf{P}_{\theta}(X_{i}\le t \ \forall i)
= \left(\frac{t}{\theta}\right)^{n},
\]
since the \(X_i\) are i.i.d.\ Uniform\((0,\theta)\).
Differentiating yields the density
\footnote{Equivalently, \(\frac{1}{\theta}T \sim \mathrm{Beta}(n,1)\).}
\[
p^{T}_{\theta}(t)
= \frac{n t^{n-1}}{\theta^{n}} \bs{1}_{[0,\theta]}(t).
\]

\begin{claim}
The estimator \(\frac{n+1}{n}T\) is the UMVUE of \(\theta\).
\end{claim}

\begin{proof}
We first show unbiasedness. Using the density of \(T\),
\[
\mathsf{E}_{\theta}[T]
= \int_{0}^{\theta} t \cdot \frac{n t^{\,n-1}}{\theta^{n}}\,dt
= \frac{n}{n+1}\theta,
\]
and hence \(\frac{n+1}{n}T\) is unbiased for \(\theta\).

To apply the Lehmann--Scheff\'e theorem, it remains to verify that \(T\) is complete.
Let \(h:\R\to\R\) satisfy
\[
\mathsf{E}_{\theta}[h(T)] = 0
\quad \text{for all } \theta>0.
\]
Then
\begin{align*}
\mathsf{E}_{\theta}[h(T)]
&= \int_{0}^{\theta} h(t)\,\frac{n t^{n-1}}{\theta^{n}}\,dt = 0
\quad \forall\,\theta>0 \\
\iff\;
& \int_{0}^{\theta} h(t)t^{n-1}\,dt = 0
\quad \forall\,\theta>0.
\end{align*}

Since the integral of \(h(t)t^{n-1}\) over every interval \([0,\theta]\) vanishes, it follows that
\[
h(t)t^{n-1}\bs{1}_{[0,\infty)}(t)=0
\quad \text{Lebesgue-a.e.}
\footnote{If a locally integrable function integrates to zero over all intervals \([0,\theta]\), then it must vanish almost everywhere; see, e.g., Shorack (2017), p.~44.}
\]
Consequently, \(h(T)=0\) \(\mathsf{P}_{\theta}\)-a.s.\ for all \(\theta>0\), and \(T\) is complete.

Since \(T\) is both sufficient and complete, the Lehmann--Scheff\'e theorem implies that
\(\frac{n+1}{n}T\) is the UMVUE of \(\theta\).
\end{proof}
\end{ex}

\section{Completeness in Exponential Families}%
\label{sec:cef}

A large and important class of models in which completeness arises almost automatically is given by exponential families.  
In this setting, completeness is closely related to uniqueness properties of Laplace and characteristic transforms; see, for instance, Lehmann and Romano (2005, Theorem~4.3.1).\footnote{The proof reduces the problem to uniqueness of characteristic functions.}

\begin{theo}[]
\label{theo:cef}
Let \(\cc{P}=\{\mathsf{P}_{\theta} : \theta\in\Theta\}\) be an exponential family with natural parameter \(\eta(\theta)\) and sufficient statistic \(T(X)\).  
That is, each \(\mathsf{P}_{\theta}\) admits a density (with respect to a \(\sigma\)--finite measure \(\nu\)) of the form
\[
p_{\theta}(x)
= \exp\!\left\{ \langle \eta(\theta), T(x) \rangle - B(\theta) \right\} h(x).
\]
If the set \(\{\eta(\theta) : \theta\in\Theta\}\) has non--empty interior, then \(T\) is complete.
\end{theo}

\paragraph{Interpretation.}
The theorem states that, for a \emph{full} exponential family, sufficiency already implies completeness.  
The condition that the natural parameter space has non--empty interior means that the model is not artificially restricted to a lower--dimensional subset of the natural parameter space.

Together with the Lehmann--Scheff\'e theorem, this result immediately yields UMVUEs in many classical models, including the binomial model (\autoref{ex:bin}).  
In contrast, the uniform model (\autoref{ex:unif}) is not an exponential family, but still admits a complete sufficient statistic via a separate argument.

\paragraph{Why varying support rules out exponential families.}
In an exponential family, the density is strictly positive wherever \(h(x)>0\), and the function \(h\) does not depend on \(\theta\).
Consequently, all distributions in the family share the same support.
Models in which the support depends on \(\theta\), such as the Uniform\((0,\theta)\) family, cannot be exponential families.

\paragraph{Role of the interior condition.}
If the natural parameter space \(\{\eta(\theta)\}\) contains a full--dimensional open set, then completeness follows.
For example, in the normal family with unknown mean \(\mu\in\R\) and variance \(\sigma^{2}>0\), the natural parameters
\[
\left(\frac{\mu}{\sigma^{2}},\,\frac{1}{\sigma^{2}}\right)
\]
range over an open subset of \(\R^{2}\), and the corresponding sufficient statistics are complete.

However, if the model is restricted, completeness may fail.
For instance, if one considers only normal distributions satisfying a constraint that links mean and variance (e.g.\ \(\mu=\sigma^{2}\)), then the natural parameters lie on a lower--dimensional subset of \(\R^{2}\).
In this case, the parameter space has empty interior, and the theorem no longer applies.

\begin{remark}[]
Completeness is a property of all measurable functions of the statistic.
Even if a statistic takes values only in \((0,\infty)\), this alone does not imply completeness, since arbitrary functions of the statistic may change signs.
Completeness requires that \emph{every} function with zero expectation for all parameter values must vanish almost surely.
\end{remark}

\medskip
Although the result may appear abstract, it is ultimately a consequence of familiar uniqueness principles for Laplace and characteristic transforms.

\begin{ex}[Estimating a Gaussian probability]
Let \(X_{1},\dots,X_{n}\) be i.i.d.\ \(\cc{N}(\mu,1)\) with unknown mean \(\mu\) and known variance, and assume \(n\ge2\).

\emph{Parameter of interest.}
For a fixed \(x\in\R\), we aim to estimate
\[
\gamma(\mu)
:= \mathsf{P}_{\mu}(X_{1}\le x)
= \Phi(x-\mu),
\]
where \(\Phi\) denotes the distribution function of \(\cc{N}(0,1)\).

\paragraph{Unbiased estimation.}
Since probabilities are expectations of indicator functions, an unbiased (but inefficient) estimator is immediately available:
\[
T(X_{1},\dots,X_{n})
= \bs{1}_{(-\infty,x]}(X_{1}),
\qquad
\mathsf{E}_{\mu}[T] = \gamma(\mu).
\]

\paragraph{Towards a UMVUE.}
The normal location model with known variance is a one--parameter exponential family.
Its sufficient statistic \(\ol{X}_{n}\) is also complete (cf.\ \autoref{sec:cef}).
Hence, by the Lehmann--Scheff\'e theorem, the UMVUE of \(\gamma(\mu)\) is obtained by conditioning \(T\) on \(\ol{X}_{n}\):
\begin{align*}
T^{*}(\bar{x}_{n})
&= \mathsf{E}_{*}\!\left[T(\bs{X}) \mid \ol{X}_{n}=\bar{x}_{n}\right] \\
&= \mathsf{P}_{*}(X_{1}\le x \mid \ol{X}_{n}=\bar{x}_{n}) \\
&= \mathsf{P}_{*}(X_{1}-\ol{X}_{n} \le x-\bar{x}_{n} \mid \ol{X}_{n}=\bar{x}_{n}).
\end{align*}

\paragraph{Evaluation.}
The pair \((X_{1}-\ol{X}_{n},\,\ol{X}_{n})\) is jointly Gaussian, being a linear transformation of \((X_{1},\dots,X_{n})\).
A direct computation shows that
\[
\mathsf{Cov}(X_{1}-\ol{X}_{n},\,\ol{X}_{n}) = 0,
\]
and hence \(X_{1}-\ol{X}_{n}\) and \(\ol{X}_{n}\) are independent.
Therefore,
\[
T^{*}(\bar{x}_{n})
= \mathsf{P}(X_{1}-\ol{X}_{n} \le x-\bar{x}_{n}).
\]

Since \(X_{1}-\ol{X}_{n} \sim \cc{N}\!\left(0,\,1-\frac{1}{n}\right)\), we obtain
    \[
        \mathsf{P}\left( \frac{X_1-\ol{X}_{n}}{\sqrt{1-\frac{1}{n}} }\le \frac{x-\bar{x}_{n}}{\sqrt{1-\frac{1}{n}} } \right) 
        = \Phi \left( (x-\bar{x}_{n}) \frac{1}{\sqrt{1-\frac{1}{n}} } \right).
    \]
\paragraph{Conclusion.}
The UMVUE of \(\gamma(\mu)=\mathsf{P}_{\mu}(X_{1}\le x)\) is
\[
T^{*}(\ol{X}_{n})
= \Phi\!\left( (x-\ol{X}_{n}) \frac{1}{\sqrt{1-\frac{1}{n}}} \right).
\]
\end{ex}

\section{Nonparametric Unbiased Estimation}%
\label{sec:nue}

Let \(X_{1}, \dots, X_{n}\) be i.i.d. random variables taking values in \(\R\) with distribution function (d.f.) \(F\).

\myparagraph{Motivation and Model}
So far we have mostly discussed parametric models, where we assume a specific distributional shape but unknown finite-dimensional parameters. For example:

\begin{itemize}
    \item \(X_i \sim \mathsf{Normal}(\mu, \sigma^2)\) with unknown \(\mu\) and \(\sigma^2\),
    \item \(X_i \sim \mathsf{Poisson}(\lambda)\) with unknown rate \(\lambda\).
\end{itemize}

In contrast, \emph{nonparametric} models are infinite-dimensional: we assume less about the shape of the distribution. For instance, we might only assume that \(F\) has a density, is smooth, or log-concave, without specifying a particular parametric family.

Here we consider a general nonparametric model. Let \(F\) belong to a class \(\cc{F}\) of distribution functions with the following properties:

\begin{itemize}
    \item \(\cc{F}\) is \emph{convex}: any convex combination of distributions in \(\cc{F}\) also belongs to \(\cc{F}\),
    \item All \(F \in \cc{F}\) are \emph{absolutely continuous} (i.e., they have a density \(f\) with respect to Lebesgue measure),
    \item \(\cc{F}\) contains all \emph{uniform distributions}.
\end{itemize}

This is a very large class of distributions: all have densities, but otherwise the assumptions are minimal. 

\subsection*{Order Statistics}
We know from general theory that for i.i.d. samples, the order of the data points contains no additional information. Therefore, the \emph{order statistics} 
\[
T(X_1, \dots, X_n) = (X_{(1)}, \dots, X_{(n)})
\] 
are \emph{sufficient} statistics. Here, \(X_{(1)}\) is the smallest observation, \(X_{(2)}\) the second smallest, and so on up to \(X_{(n)}\).

Interestingly, in this nonparametric model, the order statistics are not just sufficient—they are also \emph{complete}.

\begin{theo}[]
\label{theo:nue}
The order statistics 
\[
T(X_{1}, \dots, X_{n}) = (X_{(1)}, \dots, X_{(n)})
\]
are complete for the class \(\cc{F}\).
\end{theo}

\begin{proof}[]
To show completeness, we must prove that if \(h(T)\) satisfies 
\[
\mathsf{E}_F[h(T)] = 0 \quad \text{for all } F \in \cc{F},
\] 
then \(h(T) = 0\) almost surely for all \(F \in \cc{F}\).

\paragraph{Step 1: Consider mixture distributions.}  
Let \(F_1, \dots, F_n \in \cc{F}\) with densities \(f_1, \dots, f_n\) and weights \(\alpha_1, \dots, \alpha_n > 0\). Form the mixture
\[
F = \frac{1}{\sum_{i=1}^{n} \alpha_i} \sum_{i=1}^{n} \alpha_i F_i,
\]
which belongs to \(\cc{F}\) by convexity. Its density is 
\[
f = \frac{1}{\sum_{i=1}^{n} \alpha_i} \sum_{i=1}^{n} \alpha_i f_i.
\]

\paragraph{Step 2: Express expectation as a polynomial.}  
By assumption, \(\mathsf{E}_F[h(T)] = 0\), i.e.,
\[
\int h(T(x_1, \dots, x_n)) \prod_{j=1}^{n} f(x_j) \, dx = 0.
\]  
Substituting the mixture density gives
\[
\int h(T(x_1, \dots, x_n)) \prod_{j=1}^{n} \Big( \sum_{i=1}^{n} \alpha_i f_i(x_j) \Big) dx = 0.
\]

This is a polynomial in \(\alpha_1, \dots, \alpha_n\). Since it vanishes for all \(\alpha_i>0\), all coefficients must be zero.

\paragraph{Step 3: Focus on the coefficient of \(\prod_i \alpha_i\).}  
This coefficient is
\[
\sum_{\pi \in S_n} \int h(T(x_1, \dots, x_n)) \prod_{i=1}^{n} f_i(x_{\pi(i)}) \, dx.
\]
Reindexing the integration variables using the permutation \(\pi\) gives
\[
n! \int h(T(x_1, \dots, x_n)) \prod_{i=1}^{n} f_i(x_i) \, dx = 0.
\]

\paragraph{Step 4: Reduce to uniform distributions.}  
Choose \(f_i(x) = \frac{1}{b_i - a_i} \mathbf{1}_{[a_i, b_i]}(x)\), i.e., \(F_i \sim \text{Uniform}(a_i, b_i)\). Then
\[
0 = \int_{a_1}^{b_1} \dots \int_{a_n}^{b_n} h(T(x_1, \dots, x_n)) \, dx_1 \dots dx_n.
\]

\paragraph{Step 5: Conclude by rectangles.}  
Since this holds for all rectangles \([a_1, b_1] \times \dots \times [a_n, b_n]\), a standard result from integration theory implies
\[
h(T) \equiv 0 \quad \text{almost everywhere.}
\]

Because all \(F \in \cc{F}\) are absolutely continuous, we conclude
\[
\mathsf{P}_F(h(T) = 0) = 1 \quad \text{for all } F \in \cc{F}.
\]
\end{proof}
\paragraph{Discussion.}  
This shows that in this very general nonparametric model, the order statistics are both sufficient and complete. By the Lehmann–Scheffé theorem, any unbiased estimator of a parameter can be improved (or obtained) by conditioning on the order statistics. The completeness property here relies on the richness of \(\cc{F}\) and the convexity/mixture argument.

\subsection{U--Statistics}%
\label{sub:u-stats}

In statistics, we often want to estimate parameters that are \emph{functions of several observations at a time}, not just of single data points. This leads naturally to the notion of \textbf{U--statistics}, which generalize many familiar estimators such as the sample mean or variance.

\bigskip

A U--statistic is built from a \emph{kernel function} \(h\) that acts on \(m\) data points at a time:
\[
    h \: : \: \cc{X}^{m}\to\R, \quad \text{for }\cc{X}\subseteq \R.     
\]
We usually assume that \(h\) is \emph{symmetric}, meaning that permuting its arguments does not change its value:
\[
    h(x_{\pi(1)}, \dots, x_{\pi(m)})
    = h(x_{1}, \dots, x_{m}) \quad \forall \pi \in S_{m}, \ \forall x \in \cc{X}.
\]
If \(h\) is not symmetric, we can symmetrize it by averaging over all permutations:
\[
    \tilde{h}(x_1, \dots, x_m) = \frac{1}{m!}\sum_{\pi \in S_m} h(x_{\pi(1)}, \dots, x_{\pi(m)}).
\]

\begin{definition}[]
For \(n \ge m\), the \textbf{U--statistic} of order \(m\) with kernel \(h\) is
\[
    U_{n}(x_{1}, \dots, x_{n}) = \frac{1}{\binom{n}{m} }\sum_{1\le i_1<i_2<\dots <i_{m}\le n} h(x_{i_1}, \dots, x_{i_{m}}).
\]
Intuitively, \(U_n\) computes \(h\) on \emph{every subset of size \(m\)} of the data and averages the results. 
\end{definition}

\paragraph{Intuition}%
Think of your dataset as a vector in \(\R^n\). A U--statistic of order \(m\) examines all possible subsets of size \(m\), applies a function \(h\) to each, and averages the outputs. This approach captures interactions among \(m\) points and leads to unbiased estimators for parameters that are expectations of such functions.

\subsection{U--Statistics and UMVUE}%
\label{sub:u-stats-umvue}

We now connect U--statistics to \emph{optimal unbiased estimation}. Suppose we want to estimate a parameter that is the expectation of a kernel function \(h\) applied to \(m\) i.i.d. observations.

\begin{theo}[]
\label{theo:u-stats}
Let \(\cc{F}\) be a class of distributions satisfying:
\begin{itemize}
    \item \(\cc{F}\) is convex,
    \item \(\cc{F}\) contains only absolutely continuous distributions,
    \item \(\cc{F}\) contains all uniform distributions.
\end{itemize}
If \(X_{1}, \dots, X_{n}\) are i.i.d. with d.f. \(F\in \cc{F}\) and \(\mathsf{E}_{F}[h(X_1,\dots,X_m)^2]<\infty\) for all \(F\in \cc{F}\), then the U--statistic \(U_n\) is the \textbf{UMVUE} of
\[
    \gamma(F) = \mathsf{E}_{F}[h(X_1, \dots, X_m)].
\]
\end{theo}

\begin{proof}[]
\leavevmode
\begin{enumerate}[label=\roman*)]
    \item \textbf{Unbiasedness:} By linearity of expectation,
    \[
        \mathsf{E}_{F}[U_n(X_1,\dots,X_n)] 
        = \frac{1}{\binom{n}{m}} \sum_{1\le i_1<\dots<i_m\le n} \mathsf{E}_{F}[h(X_{i_1},\dots,X_{i_m})]
        = \gamma(F),
    \]
    because each term in the sum has expectation \(\gamma(F)\) and there are exactly \(\binom{n}{m}\) terms.
    
    \item \textbf{Optimality via Lehmann--Scheff\'e:} \(U_n\) is symmetric in its arguments, so it is a function of the \emph{order statistics} \(X_{(1)}, \dots, X_{(n)}\). Since order statistics are complete and sufficient for \(\cc{F}\), the Lehmann--Scheff\'e theorem guarantees that \(U_n\) is UMVUE.
\end{enumerate}
\end{proof}

\paragraph{Examples}%
\begin{itemize}
    \item \textbf{Sample mean:} If \(\gamma(F) = \mathsf{E}_F[X_1]\), take \(m=1\) and \(h(x_1)=x_1\). Then 
    \[
        U_n(\bs{X}) = \ol{X}_n = \frac{1}{n}\sum_{i=1}^n X_i.
    \]
    
    \item \textbf{Squared mean:} For \(\gamma(F) = (\mathsf{E}_F[X_1])^2\), use \(m=2\) and \(h(x_1,x_2) = x_1 x_2\). Then
    \[
        U_n(\bs{X}) = \frac{1}{\binom{n}{2}} \sum_{1\le i<j\le n} X_i X_j,
    \]
    which is an unbiased estimator of \((\mathsf{E}[X_1])^2\) because \(\mathsf{E}[X_i X_j] = (\mathsf{E}[X_1])^2\) for independent \(X_i, X_j\).
    
    \item \textbf{Variance:} For \(\gamma(F) = \mathsf{Var}_F[X_1]\), noting that \(\mathsf{Var}_{F}\left[X_1\right]  = \mathsf{E}_{F}\left[X_1(X_1-X_2)\right] \) and symmetrizing \(x_1(x_1-x_2)\) yields this kernel
    \[
        h(x_1,x_2) = \frac{1}{2}(x_1-x_2)^2, \quad m=2.
    \]
    The corresponding U--statistic
    \[
        U_n(\bs{X}) = \frac{1}{\binom{n}{2}} \sum_{1\le i<j \le n} \frac{(X_i-X_j)^2}{2}
    \]
    is the unbiased sample variance.
\end{itemize}

\paragraph{Takeaways}%
U--statistics provide a systematic way to construct unbiased estimators of parameters that are expectations of functions of multiple observations. The kernel \(h\) encodes which combinations of observations are relevant, and the averaging over subsets ensures unbiasedness. In large nonparametric models, U--statistics often produce the \emph{best possible unbiased estimators}, capturing quantities like the mean, variance, or higher-order moments in a unified framework.
